\section{Scope, Corpus, and Qualifications}\label{app:scope}

\paragraph{Question and reader.} This reference asks what the Catholic
tradition, as Saint Thomas Aquinas systematized it, holds about the angels:
what they are, how they know and will, how they were made and how some fell,
how they are ordered, and what they do for the world and for man; and on what
authority each part rests. It is written for a Catholic reader, student, or
teacher who wants that teaching in the order of the \emph{Summa theologiae},
with the Scriptures, the Fathers, and the saints from whom Aquinas received
it and in whom the Church has lived it.

\paragraph{Corpus and completeness.} The systematic spine is \ST{I qq.~50--64}
and \ST{I qq.~106--114}. Every article of those twenty-four questions is
treated in the body and listed once in the article census
(Appendix~\ref{app:census}); other articles of the \emph{Summa} enter where
the treatise points to them. The \emph{Celestial Hierarchy} is expounded in
all fifteen of its chapters, but not reproduced. Fathers and Doctors are
read at their own loci rather than through the \emph{Summa}'s quotation of
them; a view reported only as Aquinas reports it is attributed to Aquinas's
citation. The Greek and Latin Fathers before Dionysius, Augustine, and the
Fathers and Doctors after him receive their own expositions, and the saints
and spiritual writers of later centuries enter for the Church's lived
devotion to the angels. Later school authors and the Byzantine tradition enter as witnesses
to named questions at identified loci. Magisterial acts are those that teach
or regulate belief and devotion concerning the angels; liturgical claims name
the rite, book, calendar, and rank. Handbooks and encyclopedias served only
to locate texts.

\paragraph{Grades of authority.} Each doctrinal proposition carries one of
seven grades, which describe the standing of the proposition and not the
weight of its author:
\begin{description}[leftmargin=1.6em, style=nextline, itemsep=0.2em]
\item[Defined] taught by a conciliar definition or a profession of faith;
\item[Church teaching] taught by the ordinary magisterium, the Catechism, or a
papal act, without a definition;
\item[Common teaching] held commonly by the Fathers and theologians, without a
magisterial act that teaches it;
\item[Thomist position] Aquinas's determination where another Catholic school
determines otherwise;
\item[Disputed] an open question on which the witnesses divide without a
prevailing common teaching;
\item[Pious tradition] devotional tradition received in the Church's life,
such as the lives of the saints and the apparitions the Church honours in
her calendar; and
\item[Apocryphal] material from writings outside the canon, stated as such.
\end{description}
The grade is the compiler's judgement; its basis is the act or witness cited
beside it. Where the witnesses do not settle a grade, the lower grade is given
and the reason stands with the claim. Each question of the \emph{Summa} closes
with a frame giving the governing loci, the grades with their basis, the
authorities the articles invoke together with the further witnesses read, and
any contrary opinion.

\paragraph{Texts, translations, and terminology.} English quotations of the
\emph{Summa} follow the translation of the Fathers of the English Dominican
Province (second revised edition, 1920) as presented by New Advent; Latin key
terms follow the Leonine text as presented by the Corpus Thomisticum. The
\emph{Celestial Hierarchy} is quoted in John Parker's translation (London
and Oxford, 1899) and addressed by chapter and section. Scripture is cited
by the books, chapters, and verses of the Vulgate, with the Hebrew
numbering in brackets where it differs, and quoted from the Douay--Rheims
Bible in Challoner's revision. In its own prose the book names the persons of the Old
Testament in the Douay forms (Isaias, Ezechiel, Eliseus, Osee); quotations
keep the forms of their translators. The Fathers are quoted from the Ante-Nicene
and Nicene and Post-Nicene Fathers series where those translations exist,
and otherwise rendered by the compiler from an identified Latin or Greek
edition. Where
no public-domain English exists, as for the Latin of the councils and the
Italian of recent papal addresses, the English is the compiler's own and
the original words stand beside it in italics. Greek terms are given in
italic transliteration. Appendix~\ref{app:terminology} records the preferred
English terms with their Latin and Greek.

\paragraph{Attribution of the Dionysian corpus.} The corpus bears the name
of Dionysius the Areopagite, Saint Paul's convert at Athens (Acts 17:34), and
under that name the Church of East and West received it: Gregory the Great,
John Damascene, and Aquinas cite it as the Areopagite's. Modern scholarship
dates its composition to about the year 500. The exposition follows the text
as the Church received it and gives the modern dating only where a
historical claim depends on it; the authority of its teaching rests on that
reception.

\paragraph{Rights.} Project prose, tables, and gradings are released under
the licence stated in the colophon. Extended quotations are from
public-domain editions and translations. Official documents of the Holy See
and modern translations are quoted briefly, with attribution. Postconciliar
English liturgical texts are cited rather than reproduced. Received prayers
are given only from identified editions.

\paragraph{Limits.} The fallen angels are treated as the treatise treats
them. The rites, discipline, and pastoral practice of exorcism belong to the
study \emph{Catholic Exorcism: History, Discipline, and Pastoral Practice} in
this collection. Nothing here judges a particular case, gives a method of
deliverance, or makes any claim about a named person's spiritual state. The
teaching of other religions about spirits enters only where a Father's
argument requires it. Where the Catholic schools divide, each position is
given with the school that teaches it.

\paragraph{Review.} This reference was composed with artificial-intelligence
assistance, as the generation metadata records. No independent theological,
clerical, or ecclesiastical review has been performed, and none is claimed.
It carries no \emph{nihil obstat} or \emph{imprimatur}.
